\documentclass{article}
\usepackage[english]{babel}
\usepackage{amsmath}
\usepackage{amsthm}
\usepackage{amssymb}

\newtheorem{theorem}{Theorem}[section]
\newtheorem{lemma}[theorem]{Lemma}
\newenvironment{solution}{
\begin{proof}[Solution]}{
\end{proof}}

\begin{document}

% ============================================================
\section*{Problem 1}
Where differentiable? Analytic?

\subsubsection*{1 (a)}
\[
  f(z) = z \operatorname{Im}(z)
\]
\begin{solution}
  Write $z = x + iy$ so $\operatorname{Im}(z) = y$:
  \begin{align*}
    f(x + iy)
    &= (x + iy)y
    \\ &= xy + iy^2
  \end{align*}
  So $u = xy$ and $v = y^2$. Compute partials for Cauchy-Riemann:
  \begin{align*}
    u_x &= y & v_y &= 2y
    \\ u_y &= x & v_x &= 0
  \end{align*}
  CR requires $u_x = v_y$ and $u_y = -v_x$:
  \begin{align*}
    y = 2y &\implies y = 0
    \\ x = -0 &\implies x = 0
  \end{align*}
  The partials are continuous everywhere, so $f$ is differentiable only at
  $x = 0, y = 0$.

  Analytic nowhere, because $z = 0$ is the only point where $f$ is
  differentiable and a single point can't contain any open disk of
  differentiability.
\end{solution}

\subsubsection*{1 (b)}
\[
  f(z) = \left(\operatorname{Re}(z)\right)^2 + \left(\operatorname{Im}(z)\right)^2
\]
\begin{solution}
  Write $z$ as $x + iy$ and compute partials:
  \[
    f(x + iy) = x^2 + y^2
  \]
  So $u = x^2 + y^2$ and $v = 0$.
  \begin{align*}
    u_x &= 2x & v_y &= 0
    \\ u_y &= 2y & v_x &= 0
  \end{align*}
  $u_x = v_y \implies 2x = 0 \implies x = 0$ and
  $u_y = -v_x \implies 2y = 0 \implies y = 0$. So CR holds only at
  $x = 0, y = 0$, and since the partials are continuous everywhere, $f$ is
  differentiable there.

  Differentiable only at $z = 0$.

  Analytic nowhere, for the same reason as (a): no open disk of
  differentiability exists around $z = 0$.
\end{solution}

% ============================================================
\section*{Problem 2 (Section 2.4)}

\subsection*{2 (a) --- 2.4.8}
$f$ analytic in $D$, $\operatorname{Re} f$ or $\operatorname{Im} f$ constant
$\implies f$ constant in $D$.
\begin{proof}
  Assume $f$ is analytic in a domain $D$ and either $\operatorname{Re}(f)$ or
  $\operatorname{Im}(f)$ is constant. Show $f$ is constant.

  Write $f(z)$ via its real and imaginary parts:
  \[
    f(z) = u(x, y) + iv(x, y)
  \]
  We know either $u$ or $v$ is constant because $\operatorname{Re}(f(z)) = u$
  or $\operatorname{Im}(f(z)) = v$ is constant.

  Since $f$ is analytic in $D$, the Cauchy-Riemann equations hold everywhere in
  $D$:
  \[
    u_x = v_y \qquad u_y = -v_x
  \]

  \textbf{Case: $u$ is constant.} Then $u_x = u_y = 0$, so by CR:
  \begin{align*}
    v_y &= u_x = 0
    \\ v_x &= -u_y = 0
  \end{align*}

  \textbf{Case: $v$ is constant.} Then $v_x = v_y = 0$, so by CR:
  \begin{align*}
    u_x &= v_y = 0
    \\ u_y &= -v_x = 0
  \end{align*}

  In both cases $u_x = u_y = v_x = v_y = 0$ throughout $D$. So $u$ is constant
  in $x$ and $y$, and thus constant. Likewise $v$ is constant in $x$ and $y$,
  and thus constant.

  So $f(z) = u + iv$ must be constant.
\end{proof}

\subsection*{2 (b) --- 2.4.14}
$f$ analytic, maps $D$ onto part of a line $\implies f$ constant in $D$.
\begin{proof}
  Assume $f$ is analytic and maps all points to a line. Without loss of
  generality take the line $f(z)$ maps onto for $z \in D$:
  \[
    f(D) \subseteq L
    \quad \text{with} \quad
    L = \left\{\lambda e^{ia} + b : \lambda \in \mathbb{R}\right\}
  \]
  with constants $a \in \mathbb{R}$, $b \in \mathbb{C}$.

  If $f$ is analytic, multiplying by and adding constants remains analytic. So
  $g(z)$ is analytic:
  \begin{align*}
    g(z)
    &= f(z) \cdot e^{-ia} - be^{-ia}
    \\ &= \lambda
  \end{align*}
  As $\lambda$ is real, $g(z)$ is real valued. By Problem 10 (2.4.10, shown in
  the previous HW), $g(z)$ is constant. As $f(z) = g(z)c + d$ for constants
  $c, d$, $f$ is constant.
\end{proof}

% ============================================================
\section*{Problem 3 (Section 2.5)}

\subsection*{3 (a) --- 2.5.2}
Most general harmonic $p(x, y) = ax^2 + bxy + cy^2$.
\begin{solution}
  Take the second partials:
  \begin{align*}
    p_{xx} &= 2a
    \\ p_{yy} &= 2c
  \end{align*}
  $p$ is harmonic when $p_{xx} + p_{yy} = 0$:
  \begin{align*}
    p_{xx} + p_{yy} &= 2a + 2c = 0
    \\ \implies a &= -c
  \end{align*}
  The $bxy$ term does not appear in $p_{xx}$ or $p_{yy}$, so $b$ is
  unconstrained. So all harmonic homogeneous degree 2 polynomials can be
  written as:
  \[
    p(x, y) = ax^2 + bxy - ay^2 \quad \text{for some } a, b \in \mathbb{R}
  \]
\end{solution}

\subsection*{3 (b) --- 2.5.6}
$v$ harmonic conjugate of $u$ in $D$ $\implies uv$ harmonic in $D$.
\begin{proof}
  As $v$ is the harmonic conjugate of $u$, both must be harmonic and $f$ must
  be analytic:
  \[
    f(x, y) = u + iv \text{ is analytic}
  \]
  Because both are harmonic:
  \[
    u_{xx} + u_{yy} = 0 \qquad v_{xx} + v_{yy} = 0
  \]
  Finally, because $f$ is analytic it satisfies Cauchy-Riemann:
  \begin{align*}
    u_x &= v_y
    \\ u_y &= -v_x
  \end{align*}
  Computing the partials of $g(x, y) = uv$:
  \begin{align*}
    g_x &= u_x v + v_x u
    \\ g_{xx} &= (u_{xx} v + v_x u_x) + (v_{xx} u + u_x v_x)
    \\ &= u_{xx} v + 2 v_x u_x + v_{xx} u
  \end{align*}
  Similarly:
  \[
    g_{yy} = u_{yy} v + 2 v_y u_y + v_{yy} u
  \]
  Summing the partials:
  \begin{align*}
    g_{xx} + g_{yy}
    &= v(u_{xx} + u_{yy}) + 2(u_x v_x + u_y v_y)
    \\ &\quad + u(v_{xx} + v_{yy})
    \\ &= 0 + 0 + 2(u_x v_x + u_y v_y)
    && \text{definition of harmonic (above)}
    \\ &= 2\big(u_x v_x + (-v_x)(u_x)\big)
    && \text{Cauchy-Riemann (above)}
    \\ &= 0
  \end{align*}
  As $g_{xx} + g_{yy} = 0$, $g(x, y) = uv$ is harmonic in $D$.
\end{proof}

% ============================================================
\section*{Problem 4}
Is $u$ harmonic on $\mathbb{C}$? If so, find a harmonic conjugate.

\subsubsection*{4 (a)}
\[
  u(x, y) = 2x - x^3 + 3xy^2
\]
\begin{solution}
  Consider $x, y \in \mathbb{R}$ and take the 2nd partials of $u$:
  \begin{align*}
    u_{xx} &= -6x
    \\ u_{yy} &= 6x
  \end{align*}
  So $u_{xx} + u_{yy} = 0$ and $u$ is harmonic in the entire plane.
  \[
    f(x, y) = 2x - x^3 + 3xy^2 + iv(x, y)
  \]
  By CR:
  \begin{align*}
    u_x &= 2 - 3x^2 + 3y^2 = v_y
    \\ u_y &= 6xy = -v_x
  \end{align*}
  Integrating each:
  \begin{align*}
    v &= \int v_y \, dy = 2y - 3x^2 y + y^3 + h(x)
    \\ v &= \int v_x \, dx = \int -6xy \, dx = -3x^2 y + k(y)
  \end{align*}
  Comparing, $h(x) = C$ and $k(y) = 2y + y^3 + C$. Thus:
  \[
    v(x, y) = 2y + y^3 - 3x^2 y + C
  \]
  is the harmonic conjugate.
\end{solution}

\subsubsection*{4 (b)}
\[
  u(x, y) = 3x - x^3 + 3xy^2
\]
\begin{solution}
  Consider $x, y \in \mathbb{R}$ and take the 2nd partials of $u$:
  \begin{align*}
    u_{xx} &= -6x
    \\ u_{yy} &= 6x
  \end{align*}
  So $u_{xx} + u_{yy} = 0$ and $u$ is harmonic in the entire plane.
  \[
    f(x, y) = 3x - x^3 + 3xy^2 + iv(x, y)
  \]
  By CR:
  \begin{align*}
    u_x &= 3 - 3x^2 + 3y^2 = v_y
    \\ u_y &= 6xy = -v_x
  \end{align*}
  Integrating each:
  \begin{align*}
    v &= \int v_y \, dy = 3y - 3x^2 y + y^3 + h(x)
    \\ v &= \int v_x \, dx = \int -6xy \, dx = -3x^2 y + k(y)
  \end{align*}
  Comparing these expressions for $v$: $h(x) = C$ and $k(y) = 3y + y^3 + C$.
  Thus:
  \[
    v(x, y) = 3y + y^3 - 3x^2 y + C
  \]
  is the harmonic conjugate.
\end{solution}

\subsubsection*{4 (c)}
\[
  u(x, y) = 3x - 2x^3 + 3xy^2
\]
\begin{solution}
  Consider $x, y \in \mathbb{R}$:
  \begin{align*}
    u_{xx} &= -12x
    \\ u_{yy} &= 6x
  \end{align*}
  So $u_{xx} + u_{yy} = -6x$, which is not $0$ for $x \neq 0$. Thus $u$ is not
  harmonic in the entire plane, and it has no harmonic conjugate.
\end{solution}

% ============================================================
\section*{Problem 5}
Factor into linear/quadratic factors with real coefficients, if possible.

\subsubsection*{5 (a)}
\[
  p(z) = z^3 + z - 2
\]
\begin{solution}
  Because $p(z)$ has real coefficients it can be factored into linear and
  quadratic terms with real coefficients. We can attempt to find these roots by
  guessing and checking using polynomial long division.

  Guess: $(z - 1)$. With polynomial long division we obtain the factorization:
  \[
    p(z) = (z - 1)(z^2 + z + 2)
  \]
\end{solution}

\subsubsection*{5 (b)}
\[
  p(z) = (z - i)^2 + 2iz
\]
\begin{solution}
  Expanding $p(z)$:
  \begin{align*}
    p(z)
    &= (z - i)^2 + 2iz
    \\ &= z^2 - 2zi - 1 + 2zi
    \\ &= z^2 - 1
    \\ &= (z - 1)(z + 1)
  \end{align*}
\end{solution}

% ============================================================
\section*{Problem 6}
Taylor form centered at $z_0$.

\subsubsection*{6 (a)}
\[
  p(z) = iz^3 - 5z + 2, \quad z_0 = i
\]
\begin{solution}
  Compute each derivative and evaluate at $z_0 = i$:
  \begin{align*}
    p(z) &= iz^3 - 5z + 2 & p(i) &= i(i^3) - 5i + 2 = 1 - 5i + 2 = 3 - 5i
    \\ p'(z) &= 3iz^2 - 5 & p'(i) &= 3i(i^2) - 5 = -5 - 3i
    \\ p''(z) &= 6iz & p''(i) &= 6i(i) = -6
    \\ p'''(z) &= 6i & p'''(i) &= 6i
  \end{align*}
  All further derivatives are equal to $0$. Substituting:
  \begin{align*}
    p(z)
    &= (3 - 5i) + (-5 - 3i)(z - i) + \frac{1}{2}(-6)(z - i)^2
    + \frac{1}{6}(6i)(z - i)^3
    \\ &= (3 - 5i) + (-5 - 3i)(z - i) - 3(z - i)^2 + i(z - i)^3
  \end{align*}
  This expansion is exact because all further derivatives are equal to $0$, so
  terms with degree greater than 3 in the expansion are equal to $0$.
\end{solution}

\subsubsection*{6 (b)}
\[
  p(z) = 3i, \quad z_0 = 7 + 5i
\]
\begin{solution}
  $p(z) = 3i$ is a constant, so all derivatives are equal to zero at any point.
  Thus the Taylor expansion is just $p(z) = 3i$.
\end{solution}

\end{document}
